% !PRINTED 268
\weberpage{268}
In the group $\Lambda_p$ the normal form $S$ for the substitutions
of the third kind is contained.

If we now assume a substitution $A$ of the third kind in
the group $\Lambda_p$, then, as shall now be proved, we can transform
it into the normal form by transformation with a substitution
$U$ which itself belongs to the group $\Lambda_p$. For if we set
\weq{8}{A = \begin{pmatrix} \alpha, & \beta \\ -N\beta', & \alpha' \end{pmatrix}, \quad U = \begin{pmatrix} \lambda, & \mu \\ -N\mu', & \lambda' \end{pmatrix},}
\weq{9}{\alpha\alpha'+N\beta\beta' = 1, \qquad \lambda\lambda'+N\mu\mu' = 1,}
in which $\lambda,\lambda'$ and $\mu,\mu'$ are conjugate pairs, the
equations (3), (4), §.\,68 take the form:
\weq{10}{(\sigma-\alpha)(\sigma-\alpha')+N\beta\beta' = \sigma^2-\sigma(\alpha+\alpha')+1 = 0,}
\weq{11}{N\,\dfrac{\mu'}{\lambda} = \dfrac{\alpha-\sigma}{\beta}, \qquad \dfrac{\lambda'}{\mu} = -\dfrac{\alpha-\sigma^{-1}}{\beta},}
and the equation (10), since $A$ is to be of the third kind,
must have two roots reciprocal to one another and conjugate
imaginary.

But from the second equation (11) it follows, by passing
to the conjugate imaginary numbers, that
\[
\dfrac{\lambda}{\mu'} = -\dfrac{\alpha'-\sigma}{\beta'},
\]
and by (10) this is a consequence of the first equation (11).
If we therefore set
\[
N\mu' = \varkappa\,(\alpha-\sigma), \qquad N\mu = \varkappa'\,(\alpha'-\sigma^{-1})
\]
\[
\lambda = \varkappa\beta, \qquad \lambda' = \varkappa'\beta',
\]
in which $\varkappa,\varkappa'$ are two conjugate imaginary numbers, the
equations (11) are satisfied, and one can then determine $\varkappa\varkappa'$
so that $\lambda\lambda'+N\mu\mu' = 1$.

From this there follows yet another important result. If we
denote by $\gamma$ a primitive root of the field $\mathfrak{E}_{2,p}$, then for
the substitutions of the normal form for the second and
third kind we may assume the expressions (§.\,68):
\[
S_2^r = \begin{pmatrix} \gamma^{(p+1)r}, & 0 \\ 0, & \gamma^{-(p+1)r} \end{pmatrix}, \qquad S_3^r = \begin{pmatrix} \gamma^{(p-1)r}, & 0 \\ 0, & \gamma^{-(p-1)r} \end{pmatrix}.
\]

If now $A$ is a substitution of the second kind from $L_p$ and $A$ a
%CONT>>
